\documentclass[10pt,landscape]{article}
\usepackage{multicol}
\usepackage{calc}
\usepackage{ifthen}
\usepackage[landscape]{geometry}
\usepackage{xcolor}
\usepackage{listings}

\ifthenelse{\lengthtest { \paperwidth = 11in}}
	{ \geometry{top=.5in,left=.5in,right=.5in,bottom=.5in} }
	{\ifthenelse{ \lengthtest{ \paperwidth = 297mm}}
		{\geometry{top=1cm,left=1cm,right=1cm,bottom=1cm} }
		{\geometry{top=1cm,left=1cm,right=1cm,bottom=1cm} }
	}

\pagestyle{empty}

% Section title color
\definecolor{sblue}{HTML}{1a5276}

% Soft section background palette
\definecolor{secred}{HTML}{f5b7b1}      % triage / kick off
\definecolor{secorange}{HTML}{f5cba7}   % firewall
\definecolor{secpurple}{HTML}{d7bde2}   % users / accounts
\definecolor{secyellow}{HTML}{fdebd0}   % persistence
\definecolor{secgreen}{HTML}{abebc6}    % logs / IOCs
\definecolor{secteal}{HTML}{a3e4d7}     % files / process / network
\definecolor{secgray}{HTML}{d5dbdb}     % hardening / IR

% Code block colors
\definecolor{codebg}{HTML}{f4f6f7}
\definecolor{codeborder}{HTML}{d5dbdb}
\definecolor{codegreen}{HTML}{27ae60}
\definecolor{codekw}{HTML}{2980b9}
\definecolor{codegray}{HTML}{7f8c8d}

% Listings style
\lstset{
  basicstyle=\ttfamily\footnotesize,
  backgroundcolor=\color{codebg},
  frame=single,
  rulecolor=\color{codeborder},
  framesep=2pt,
  xleftmargin=3pt,
  xrightmargin=3pt,
  aboveskip=4pt,
  belowskip=2pt,
  breaklines=true,
  breakatwhitespace=false,
  commentstyle=\color{codegreen},
  keywordstyle=\color{codekw}\bfseries,
  stringstyle=\color{codegray},
  showstringspaces=false,
  tabsize=2,
  language=bash,
  morekeywords={iptables,systemctl,journalctl,usermod,gpasswd,visudo,pkill,killall,fuser,lsof,tcpdump,modprobe},
  morecomment=[l]{\#},
  literate={<}{{\textless}}1 {>}{{\textgreater}}1,
}

% Custom colored section: \colorsec{bgcolor}{Title}
\newcommand{\colorsec}[2]{%
  \vspace{1.5ex plus 0.5ex minus 0.5ex}%
  \noindent\rule{\linewidth}{0.4pt}\\[-0.7ex]%
  \noindent\colorbox{#1}{\parbox{\dimexpr\linewidth-2\fboxsep}{%
    \strut\textbf{\normalsize\color{sblue} #2}}}%
  \vspace{0.5ex}\par\nobreak%
}

\makeatletter
\renewcommand{\subsection}{\@startsection{subsection}{2}{0mm}%
                                {-1explus -.5ex minus -.2ex}%
                                {0.5ex plus .2ex}%
                                {\normalfont\normalsize\bfseries\color{sblue}}}
\makeatother

\setcounter{secnumdepth}{0}

\setlength{\parindent}{0pt}
\setlength{\parskip}{0pt plus 0.5ex}

% -----------------------------------------------------------------------

\begin{document}

\raggedright
\footnotesize
\begin{multicols}{3}

\setlength{\premulticols}{1pt}
\setlength{\postmulticols}{1pt}
\setlength{\multicolsep}{1pt}
\setlength{\columnsep}{2pt}

\begin{center}
     \Large{\textbf{Linux Blue Team Cheat Sheet --- CCDC}} \\
\end{center}

\colorsec{secred}{Initial Triage}
\begin{lstlisting}
w                       # who is logged in
who                     # current sessions
last -n 20              # recent logins
lastb -n 20             # failed logins
ps auxf                 # full process tree
ss -tulnp               # listening ports + PIDs
ss -anp                 # all connections + PIDs
netstat -tulnp          # alt: listening sockets
\end{lstlisting}

\colorsec{secred}{Kick Red Teamers Off}

\subsection{Kill User Sessions}
\begin{lstlisting}
# Find their PTS/TTY from `w` output
pkill -KILL -t pts/2      # kill by terminal
pkill -KILL -u baduser    # kill all by user
\end{lstlisting}

\subsection{Kill SSH Sessions}
\begin{lstlisting}
ps aux | grep sshd        # find sshd child
kill -9 <PID>
# Block their IP immediately
iptables -A INPUT -s <IP> -j DROP
\end{lstlisting}

\subsection{Kill Reverse Shells / Meterpreter}
\begin{lstlisting}
# Find suspicious outbound connections
ss -anp | grep ESTAB
lsof -i -P -n | grep ESTAB
# Look for: bash/sh/python/perl/nc/ncat
# with connections to odd ports
kill -9 <PID>
# Block C2 IP
iptables -A OUTPUT -d <C2_IP> -j DROP
\end{lstlisting}

\subsection{Kill by Port}
\begin{lstlisting}
fuser <port>/tcp          # find PID on port
lsof -i :<port>           # detailed info
kill -9 <PID>
\end{lstlisting}

\subsection{Kill Signals Quick Ref}
\begin{lstlisting}
kill -9 <PID>             # SIGKILL (force)
kill -15 <PID>            # SIGTERM (graceful)
killall <name>            # by process name
pkill -f "pattern"        # by cmdline match
\end{lstlisting}

\colorsec{secorange}{iptables Firewall}
\begin{lstlisting}
iptables -L -n -v --line-numbers

# Allow needed services
iptables -A INPUT -p tcp --dport 22 -j ACCEPT
iptables -A INPUT -p tcp --dport 80 -j ACCEPT
iptables -A INPUT -p tcp --dport 443 -j ACCEPT

# Allow established + loopback
iptables -A INPUT -m state \
  --state ESTABLISHED,RELATED -j ACCEPT
iptables -A INPUT -i lo -j ACCEPT

# Default deny
iptables -P INPUT DROP
iptables -P FORWARD DROP

# Block specific attacker IP
iptables -I INPUT -s <IP> -j DROP
iptables -I OUTPUT -d <IP> -j DROP

# Delete rule by line number
iptables -D INPUT <line#>

# Persist / restore rules
iptables-save > /etc/iptables.rules
iptables-restore < /etc/iptables.rules
\end{lstlisting}

\colorsec{secpurple}{User \& Account Hardening}
\begin{lstlisting}
# List all users with login shells
grep -v '/nologin\|/false' /etc/passwd

# Check for UID 0 (root-equiv) accounts
awk -F: '$3==0 {print}' /etc/passwd

# Lock suspicious account
usermod -L suspectuser
passwd -l suspectuser

# Change password
passwd <user>

# Check / edit sudoers
visudo
ls -la /etc/sudoers.d/

# Remove user from sudo/wheel
gpasswd -d baduser sudo
gpasswd -d baduser wheel

# Find SSH key backdoors
find / -name authorized_keys 2>/dev/null
# Review & remove unknown keys
\end{lstlisting}

\colorsec{secyellow}{Hunting Persistence}

\subsection{Cron Jobs}
\begin{lstlisting}
crontab -l               # current user
crontab -l -u root       # root crontab
ls -la /etc/cron*         # system crons
cat /etc/crontab
# Check all users
for u in $(cut -f1 -d: /etc/passwd); do
  crontab -l -u $u 2>/dev/null; done
# Also: /var/spool/cron/crontabs/
\end{lstlisting}

\subsection{Systemd Services / Timers}
\begin{lstlisting}
systemctl list-units --type=service \
  --state=running
systemctl list-timers --all
# Look for unfamiliar unit files in:
ls /etc/systemd/system/
ls /usr/lib/systemd/system/
ls ~/.config/systemd/user/
# Disable rogue service
systemctl stop <svc>
systemctl disable <svc>
\end{lstlisting}

\subsection{Init / Startup Scripts}
\begin{lstlisting}
ls -la /etc/init.d/
cat /etc/rc.local
\end{lstlisting}

\subsection{Shell Profiles (Backdoor Triggers)}
\begin{lstlisting}
cat /etc/profile /etc/bash.bashrc
cat ~/.bashrc ~/.bash_profile ~/.profile
# Look for: curl | wget | nc | /dev/tcp
# Also check: /etc/environment
#   /etc/ld.so.preload (LD_PRELOAD hijack)
\end{lstlisting}

\subsection{SUID / SGID Binaries}
\begin{lstlisting}
find / -perm -4000 -type f 2>/dev/null
find / -perm -2000 -type f 2>/dev/null
# Compare against known-good baseline
\end{lstlisting}

\subsection{Kernel Modules}
\begin{lstlisting}
lsmod                    # list loaded modules
modprobe -r <module>     # remove a module
# Look for unknown / rootkit modules
\end{lstlisting}

\colorsec{secgreen}{Log Analysis --- IOCs}

\subsection{Key Log Locations}
\begin{tabular}{@{}ll@{}}
\verb!/var/log/auth.log!   & Auth (Debian/Ubuntu) \\
\verb!/var/log/secure!     & Auth (RHEL/CentOS) \\
\verb!/var/log/syslog!     & General (Debian) \\
\verb!/var/log/messages!   & General (RHEL) \\
\verb!/var/log/apache2/!   & Apache web logs \\
\verb!/var/log/nginx/!     & Nginx web logs \\
\verb!/var/log/cron!       & Cron execution \\
\verb!/var/log/kern.log!   & Kernel messages \\
\verb!journalctl -xe!      & Systemd journal \\
\end{tabular}

\subsection{IOC Searches}
\begin{lstlisting}
# Brute force SSH
grep "Failed password" /var/log/auth.log
# Count IPs with most failures
grep "Failed password" /var/log/auth.log \
  | awk '{print $(NF-3)}' | sort | \
  uniq -c | sort -rn | head

# Successful logins
grep "Accepted" /var/log/auth.log

# Privilege escalation
grep "sudo:" /var/log/auth.log
grep "su:" /var/log/auth.log

# New user creation
grep "useradd\|adduser" /var/log/auth.log

# Password changes
grep "passwd" /var/log/auth.log

# Web shells in access logs
grep -i "cmd=\|exec\|system\|passthru\
\|eval\|shell" /var/log/apache2/access.log

# Log tampering: check timestamps/sizes
ls -la /var/log/
\end{lstlisting}

\colorsec{secteal}{Suspicious Files}
\begin{lstlisting}
# Recently modified files (last 10 min)
find / -mmin -10 -type f 2>/dev/null

# World-writable files
find / -perm -0002 -type f 2>/dev/null

# Attacker drop zones
ls -la /tmp /dev/shm /var/tmp

# Hidden files in odd places
find / -name ".*" -type f 2>/dev/null \
  | grep -v -e proc -e sys

# Recent webshells
find /var/www -name "*.php" -mmin -60

# Files owned by nobody/nogroup
find / -nouser -o -nogroup 2>/dev/null

# Executables in /tmp (should not exist)
find /tmp -type f -executable
\end{lstlisting}

\colorsec{secteal}{Process Forensics}
\begin{lstlisting}
ps auxww                 # full command lines

# Inspect a specific process
ls -la /proc/<PID>/exe   # actual binary
cat /proc/<PID>/cmdline  # command line
ls -la /proc/<PID>/fd/   # open file handles
cat /proc/<PID>/environ  # environment vars

# DELETED but running binaries (IOC!)
ls -la /proc/*/exe 2>/dev/null | grep deleted

# Processes running from /tmp (IOC!)
ls -la /proc/*/exe 2>/dev/null | grep tmp
\end{lstlisting}

\colorsec{secteal}{Network Forensics}
\begin{lstlisting}
# Active connections with PIDs
ss -anp | grep ESTAB

# ARP table (look for spoofing)
arp -a
ip neigh

# Routing table (check for hijacks)
ip route
route -n

# Promiscuous mode (sniffers running?)
ip link | grep PROMISC

# Capture traffic for analysis
tcpdump -i eth0 -w capture.pcap
tcpdump -i eth0 port not 22

# DNS (check for hijacking)
cat /etc/resolv.conf
\end{lstlisting}

\colorsec{secgray}{SSH Hardening}
Edit \verb!/etc/ssh/sshd_config!:
\begin{lstlisting}
PermitRootLogin no
PasswordAuthentication no
MaxAuthTries 3
AllowUsers user1 user2
Protocol 2
\end{lstlisting}
Then: \verb!systemctl restart sshd!

\colorsec{secgray}{Important Config Files}
\begin{tabular}{@{}ll@{}}
\verb!/etc/passwd!          & User accounts \\
\verb!/etc/shadow!          & Password hashes \\
\verb!/etc/group!           & Groups \\
\verb!/etc/sudoers!         & Sudo privileges \\
\verb!/etc/ssh/sshd_config! & SSH server config \\
\verb!/etc/hosts!           & Static DNS entries \\
\verb!/etc/resolv.conf!     & DNS servers \\
\verb!/etc/crontab!         & System cron \\
\verb!/etc/fstab!           & Mount points \\
\verb!/etc/exports!         & NFS shares \\
\verb!/etc/samba/smb.conf!  & Samba config \\
\verb!/etc/apache2/!        & Apache config \\
\verb!/etc/nginx/!          & Nginx config \\
\verb!/etc/mysql/my.cnf!    & MySQL config \\
\end{tabular}

\colorsec{secgreen}{What to Hunt For (IOC Names)}
\textbf{Suspicious process names:}
\begin{itemize} \itemsep0pt \parskip0pt \parsep0pt \topsep0pt
\item Netcat variants: nc, ncat, netcat, ncat.traditional
\item Tunneling: chisel, ligolo, socat, ngrok, frp
\item Enum scripts: linpeas, linenum, pspy, lse.sh
\item Credential dumpers: mimipenguin, lazagne, keydump
\item Scanners: nmap, masscan, rustscan, zmap
\item Stagers: msfvenom output, elf with random names in /tmp
\end{itemize}

\textbf{Suspicious behaviors:}
\begin{itemize} \itemsep0pt \parskip0pt \parsep0pt \topsep0pt
\item bash/sh/python with open network socket
\item Process running from /tmp, /dev/shm, /var/tmp
\item Deleted binary still running (see /proc/*/exe)
\item Outbound connection on high or unusual port
\item Cron / systemd launching unknown binary
\item \verb!/etc/ld.so.preload! non-empty (LD\_PRELOAD hook)
\item \verb!authorized_keys! added to unexpected users
\end{itemize}

\colorsec{secgray}{Protect Key Files (chattr)}
\begin{lstlisting}
# Make file immutable -- even root cannot modify
chattr +i /etc/passwd /etc/shadow /etc/sudoers
chattr +i /etc/hosts /etc/resolv.conf
chattr +i /etc/ssh/sshd_config
# Remove immutable flag when you need to edit
chattr -i /etc/passwd
# List file attributes
lsattr /etc/passwd
\end{lstlisting}

\colorsec{secteal}{Auditd --- File \& Auth Monitoring}
\begin{lstlisting}
systemctl start auditd
# Watch critical files for writes/attr changes
auditctl -w /etc/passwd  -p wa -k passwd_mod
auditctl -w /etc/shadow  -p wa -k shadow_mod
auditctl -w /etc/sudoers -p wa -k sudoers_mod
auditctl -w /etc/ssh/sshd_config -p wa -k ssh_mod
# Watch for privilege escalation attempts
auditctl -w /usr/bin/sudo -p x -k sudo_exec
# Query logs
ausearch -k passwd_mod
ausearch -k sudo_exec | grep -i "res=failed"
aureport --auth --summary
\end{lstlisting}

\colorsec{secorange}{TCP Wrappers (hosts.allow/deny)}
\begin{lstlisting}
# Deny all by default
echo "ALL: ALL" >> /etc/hosts.deny
# Then whitelist what you need in hosts.allow
echo "sshd: 192.168.1.0/24" >> /etc/hosts.allow
echo "sshd: 10.0.0.0/8"     >> /etc/hosts.allow
# Verify wrapper is active (check sshd linkage)
ldd /usr/sbin/sshd | grep libwrap
\end{lstlisting}

\colorsec{secpurple}{Account Lockout \& Password Policy}
\begin{lstlisting}
# Force immediate password expiry (must reset)
chage -d 0 <user>
# Set min/max password age
chage -m 1 -M 90 <user>
# View expiry info
chage -l <user>
# Lock account after N failed logins (PAM)
# /etc/pam.d/common-auth (Debian) or system-auth
# Add: auth required pam_tally2.so deny=5
#       account required pam_tally2.so
# Reset counter: pam_tally2 --user=<user> --reset
# Expire ALL regular-user passwords at once
for u in $(awk -F: '$3>=1000{print $1}' /etc/passwd); do
  chage -d 0 $u; done
\end{lstlisting}

\colorsec{secgray}{Hardening Checklist}
\begin{itemize} \itemsep0pt \parskip0pt \parsep0pt \topsep0pt
\item Change ALL default passwords immediately
\item \verb!PermitRootLogin no! in sshd\_config
\item \verb!PasswordAuthentication no! + key-based auth
\item Disable unused services: \verb!systemctl disable <svc>!
\item Update packages: \verb!apt update \&\& apt upgrade -y!
\item \verb!/etc/shadow! permissions: \verb!chmod 640!, owner root:shadow
\item Remove unnecessary SUID bits: \verb!chmod u-s <bin>!
\item \verb!/etc/hosts.deny!: ALL: ALL (TCP wrappers)
\item \verb!chattr +i! on /etc/passwd, shadow, sudoers
\item Disable IPv6 if not needed: \verb!sysctl net.ipv6.conf.all.disable\_ipv6=1!
\item Check \verb!/etc/hosts! for malicious redirects
\item Verify DNS servers in \verb!/etc/resolv.conf!
\item Review \verb!/etc/ld.so.preload! --- should be empty
\item Back up configs before making changes
\end{itemize}

\colorsec{secgray}{Incident Response Flow}
\begin{enumerate} \itemsep0pt \parskip0pt \parsep0pt \topsep0pt
\item \textbf{Identify:} \verb!w!, \verb!ss -anp!, \verb!ps auxf!
\item \textbf{Kill:} sessions, reverse shells, C2 processes
\item \textbf{Block:} \verb!iptables! attacker IPs; TCP wrappers
\item \textbf{Lockdown:} passwords, accounts, SSH keys, sudoers
\item \textbf{Protect:} \verb!chattr +i! on critical files
\item \textbf{Hunt:} cron, systemd, shell profiles, SUID, ld.so.preload
\item \textbf{Harden:} firewall default-deny, disable unused services
\item \textbf{Monitor:} auditd, watch logs \& active connections
\end{enumerate}

\rule{0.3\linewidth}{0.25pt}
\scriptsize

CCDC Blue Team --- Linux | \today

\end{multicols}
\end{document}
